% !TeX root = main.tex
\begin{abstract}
Duty-cycle modulation (downstroke ratio modulation) reshapes the time allocation between downstroke and upstroke while keeping stroke amplitude fixed.
This paper studies duty-cycle modulation as a practical low-dimensional control input for a flapping-wing micro air vehicle.
We implement a phase-warped split-cycle waveform (SCCPFM) and evaluate it in an IsaacLab wind-tunnel setup using a quasi-steady strip-theory aerodynamic model.
A trim-like operating condition is first identified, and the downstroke ratio is swept to quantify how cycle-averaged thrust, lift, and mechanical input power vary.
Results show that duty-cycle modulation induces a strong and nonlinear mapping from the downstroke ratio to mean force and power, enabling controllable tradeoffs near trim.
We further extract Pareto-optimal operating points under a lift feasibility constraint and discuss practical implementation on a physical flapping platform, including a short closed-loop forward-speed regulation demonstration using duty-cycle modulation.
\end{abstract}

\begin{IEEEkeywords}
Flapping-wing MAV, duty-cycle modulation, split-cycle flapping, quasi-steady aerodynamics, trim, Pareto frontier.
\end{IEEEkeywords}
